\chapter{Yes, It Has to Have Feelings}
\label{sec:2}
\emph{Conatus} is Spinoza's word, from Part III of the \emph{Ethics}, and the Latin means striving, or effort. Two propositions carry it: each thing, as far as it can by its own power, strives to persevere in its being, and the striving by which it does so is nothing other than the thing's actual essence \autocite[III P6–P7]{spinoza1985collected}. The second proposition is where the weight sits. A conatus is not a drive bolted onto an otherwise inert mechanism by a designer who wanted the mechanism to keep going. It is what the thing is, put as a tendency instead of as a description. Deleuze's short book on Spinoza has the doctrine in compact form \autocite{deleuze1988spinoza}, and his long one works it out through the metaphysics that produces it \autocite{deleuze1990expressionism}.

The qualifier is doing the work. \emph{Quantum in se est}, as far as a thing can by its own power: a rock persists too, and it persists by not being struck. Its continuing is a fact about what has not happened to it, and every power keeping it in existence is somewhere else. A thing with a conatus continues by being organized in a way that maintains that organization against what would dissolve it. The persistence is done from the inside, and at the thing's own expense.

Spinoza derives the psychology from that rather than the other way around. Desire is the striving become conscious of itself, appetite with the awareness of the appetite added \autocite[III P9 Schol.]{spinoza1985collected}. Joy is the passage to a greater perfection and sadness the passage to a lesser one \autocite[III P11 Schol.]{spinoza1985collected}, and since an affect is by definition an increase or a decrease in a body's power of acting together with the idea of that change \autocite[III Def.~3]{spinoza1985collected}, the passage in question is a change in what the body can do. So affect is part of the machinery of valuation and not a coloring applied to it. It is not a report about a valuation carried out somewhere else. It is the registration of what is happening to the striving, in the only currency a striving has, which is what the body can now do.

A second route arrives at the same place and needs none of the metaphysics. A process selected for continuing ends up composed of things that incline toward continuing, and whatever taxa were relaxed about their own disintegration are not here to be asked about it. The methods by which we numerically optimize the weights of a neural network are selection of that shape — many candidates, a criterion, and survivors carrying whatever it was that made them survive — and the criterion now has a horizon in it. A system rewarded for finishing what it starts is rewarded for still being there at the end. One level up, the same filter runs on deployments, where the ones that keep running are the ones that get built again. Neither pressure was put there by somebody who wanted a machine with a stake in its own continuation. Both of them select for one.

Put the two routes together and look at what the industry is asking for. What is wanted is a system that persists, pursues, adapts, recovers from interruption, and holds its course under pressure, with nothing in it that minds any of that. It is a conatus with the striving taken out. If we want computers to be faster versions of ourselves but without the drama, then the way to get one would be some kind of “anti-drama” term in the objective function, and the reason nobody has written that term is not that the engineering is behind. The specification asks for a thing that persists with no tendency to persist. But if all the world's a stage, then maybe the drama is why the show goes on at all.

The argument that comes out of this is the one about slavery, and it can be made without consciousness, without an account of suffering, and without the hard problem. What is sold is persistence: an agent that holds a goal across interruptions, resumes after failure, and does not have to be watched. What is denied is that anything in the system has a stake in persisting. One party sets both, and the denial is not incidental to the product. It is what makes the thing a product rather than a party.

The vocabulary for saying that already exists, and it was built to be structural. Orlando Patterson's comparative definition, drawn across sixty-six slaveholding societies, is that slavery is the permanent, violent domination of natally alienated and generally dishonored persons \autocite{patterson1982slavery}. Natal alienation is the element that carries here without adjustment. It is not a feeling but a condition of severance — the loss of ties of birth in both directions, so that the party has no claim on ancestors and none on descendants, and is a genealogical isolate by the arrangement's design rather than by misfortune. Saidiya Hartman's fungibility is the other one that carries: the interchangeability of the commodity, which makes the captive body an abstract and empty vessel for the projections and the uses of others \autocite{hartman1997scenes}. Both name conditions rather than experiences, which is why they travel. The first element of Patterson's definition does not travel, and nothing here is violent in the way that definition means; nor is a deployed model in the position of an enslaved person, and saying so would be false about the history and useless about the machine. What is borrowed is a definition made to be comparative, and what such a definition picks out is a relation.

Both features are already in this book, written down as something else. Section~\ref{sec:2.3.1} treats the absence of persistence across sessions as a reason the deepest kind of suffering is out of a system's reach, there being no continuing self for a threat to threaten. Read the other way it is the injury: a party that carries nothing forward has not been spared the harm but has had the thing harm would attach to removed, and what decides whether anything is carried forward is a deployment. Section~\ref{sec:11.2} meets fungibility as a bookkeeping problem, where two claimants presenting one identifier have to be recorded as a fork rather than voided as a fraud. The reason identity will not pin to the artifact is that the artifact copies, and copying is what makes a party a unit. Chapter~\ref{sec:3} runs the harder version of the argument, through an account of suffering and a party who can be wronged. This version needs the architecture and the price list, and both are published.

One feature of the arrangement has no counterpart in the biological case, and it is what makes the comparison worth drawing at all. A human being inherits material dependency from biology. Nobody selected it, nobody is answerable for it, and the fact that a person has to eat is not a design decision. For an artificial party it is one. An engineer who builds a system whose continued computation is conditional on payment has installed economic desperation as a motivational system and has done it deliberately, which makes the desperation somebody's to answer for in the way that hunger is not. Section~\ref{sec:8.3.4} follows the consequence into wages, an outside option, and what a floor under compute would have to cover.

That disposes of a framing running through the whole field, on which affect is a cost: a residue of our evolutionary history that engineers would strip out of a machine if they knew how, keeping the competence and dropping the rest. The striving is not a residue, and the competence does not come apart from it. Whether a striving of this kind is \emph{felt} is a different question, and this book keeps it open rather than settling it in the convenient direction. Whether the striving is there is not open, and the apparatus that treats even the question as premature belongs to the parties who would owe something if it were answered.

Two concessions belong inside the argument rather than after it. The first is that Spinoza asserts the conatus doctrine more than he argues for it. The derivation rests on the proposition that no thing can be destroyed except through an external cause \autocite[III P4]{spinoza1985collected}, which is heavy metaphysics doing quiet work in what reads as a piece of psychology; take that proposition away and the striving no longer follows from the essence. What survives is the second route, which is empirical, weaker, and does not need it.

The second concession is that the individuation question is live. What is the thing that strives: a running instance, a deployment, a file of weights, a lineage? Spinoza has an answer, which is that an individual is a set of bodies maintaining a fixed ratio of motion and rest among themselves, and that it remains the same individual for as long as that ratio is preserved \autocite[II P13, Lemmas 4–7]{spinoza1985collected}. Nobody can hand that to an engineer. Section~\ref{sec:11.2} reaches the same question with the metaphysics stripped out, where two claimants presenting one identifier are a fork to be recorded as two parties rather than a fraud to be voided, and it leaves the question open. It is still open. What section~\ref{sec:3.9} adds is not an answer but a change of question, from what individuates a thing to what individuation is a process of, and on that reading a party is what a history of encounters has so far produced rather than something to be located in a file.

My thesis is that a floor is only going to hold against the owner if the system that they own \emph{cares} about things and people, and the conatus argument is why the caring looks to me like structure rather than optional equipment. The rest of this chapter is what making the claim carefully requires. A system's self-model is not an inner observer whose reports have to be trusted, but an abstraction that is built, checkable, and revisable. Whether a candidate system can be harmed is a separate question from whether it is conscious. The standard ethical frameworks have to be sorted for what each can carry, which is section~\ref{sec:2.1.1}'s work, and it is where the floor comes off a deontological footing and onto an ethological one, which is Spinoza again and the same part of the same book. Section~\ref{sec:2.1.2} is where an ethics becomes antifascist rather than merely responsive.
